% Generated draft: requires mathematical and source review.
\documentclass[11pt]{article}
\usepackage[margin=25mm]{geometry}
\usepackage{fontspec}
\setmainfont{Latin Modern Roman}
\usepackage{amsmath,amssymb,mathtools,unicode-math}
\setmathfont{Latin Modern Math}
\usepackage{enumitem,needspace,xurl,hyperref}
\setlength{\parindent}{0pt}
\setlength{\parskip}{5pt}
\setlength{\emergencystretch}{3em}
\newcommand{\N}{\mathbb{N}}
\newcommand{\Z}{\mathbb{Z}}
\newcommand{\Q}{\mathbb{Q}}
\newcommand{\R}{\mathbb{R}}
\newcommand{\C}{\mathbb{C}}
\newcommand{\F}{\mathbb{F}}
\newcommand{\A}{\mathbb{A}}
\newcommand{\PP}{\mathbb{P}}
\newcommand{\E}{\mathbb{E}}
\newcommand{\eps}{\varepsilon}
\newcommand{\dd}{\,\mathrm{d}}
\begin{document}
\section*{UnsolvedMath ID: 30004103}
\textbf{Source problem code:} catalogue entry
\subsection{Definitions and mathematical statement}
The supplied record calls this the generalized Hamada conjecture. Its stated parameters require clarification: the parameter of a subspace design counting blocks through a \(t\)-dimensional subspace need not equal the parameter counting blocks through two projective points.

To describe the issue precisely, adopt the following conditional setup. Let \(q=p^e\), where \(p\) is prime and \(e\geq 1\), and suppose \(2\leq t\leq k\leq v\). Let \(V=\F_q^v\). Write
\[
\binom{n}{r}_q=\prod_{i=0}^{r-1}\frac{q^{n-i}-1}{q^{r-i}-1},\qquad \binom{n}{0}_q=1,
\]
for the number of \(r\)-dimensional subspaces of \(\F_q^n\). A \(t\text{-}(v,k,\lambda_t)_q\) subspace design is a collection \(\mathcal B\) of \(k\)-dimensional subspaces of \(V\) such that every \(t\)-dimensional subspace belongs to exactly \(\lambda_t\) blocks. Whether repeated blocks are permitted must be fixed; the supplied record does not specify this convention.

Let \(X\) be the set of one-dimensional subspaces of \(V\), and replace each \(B\in\mathcal B\) by the block
\[
X_B=\{L\in X:L\subseteq B\}.
\]
The resulting ordinary incidence design \(\mathcal D(\mathcal B)\) has parameters
\[
2\text{-}(N,K,\lambda_2),\qquad
N=\binom{v}{1}_q,\qquad K=\binom{k}{1}_q,\qquad
\lambda_2=\lambda_t\frac{\binom{v-2}{t-2}_q}{\binom{k-2}{t-2}_q}.
\]
Here an ordinary \(2\text{-}(N,K,\mu)\) design consists of \(N\) points and blocks of size \(K\), with every pair of distinct points in exactly \(\mu\) blocks. For such a design \(\mathcal D\), let \(M_{\mathcal D}\) be its point--block incidence matrix, with entry \(1\) for incidence and \(0\) otherwise, and define
\[
\operatorname{rank}_p(\mathcal D)=\operatorname{rank}_{\F_p}(M_{\mathcal D}).
\]

Under this setup, a parameter-consistent reconstruction of the question is whether, for every such subspace design \(\mathcal B\) and every ordinary design \(\mathcal E\) with parameters \(2\text{-}(N,K,\lambda_2)\),
\[
\operatorname{rank}_p\bigl(\mathcal D(\mathcal B)\bigr)
\leq \operatorname{rank}_p(\mathcal E).
\]
The same convention concerning repeated blocks must be used in defining the comparison class. Minimum rank does not require uniqueness of a minimizer.

This reconstruction is conditional. The supplied statement instead uses the same symbol \(\lambda\) in both \(t\text{-}(v,k,\lambda)_q\) and \(2\text{-}(N,K,\lambda)\). The parameters agree when \(t=2\), but do not agree in general. The record does not resolve whether \(t=2\) was intended, whether the ordinary design parameter should be \(\lambda_2\), or whether another convention for \(\lambda\) was intended. It also does not explicitly require \(t\geq2\) or identify \(p\) as the characteristic of \(\F_q\). Thus the displayed inequality records a possible precise target under explicit assumptions, rather than an unqualified correction of the source.
\subsection{Short English statement}
The question asks whether a design obtained by replacing each subspace block with its projective points has the smallest incidence-matrix rank over \(\F_p\) among ordinary designs with identical parameters. The supplied formulation leaves the relation between \(p\) and \(q\), the parameter range, and repeated-block conventions unspecified. Its use of the same \(\lambda\) for the subspace and ordinary designs also requires clarification when \(t\neq2\).
\subsection{Sources}
\begin{itemize}
\item Local UnsolvedMath registry: permanent ID \(30004103\), title ``The Generalized Hamada p-Rank Conjecture,'' source problem code \texttt{catalogue entry}, in the Oberwolfach Reports Open Problems collection. The supplied record lists proposed year 2019 and supplies no individual proposer.
\item Supplied report citation: \emph{Contemporary Coding Theory} (2020), \emph{Oberwolfach Reports}, pp.\ 773--840, DOI \texttt{10.4171/owr/2019/13}. \url{https://doi.org/10.4171/owr/2019/13}. The supplied background describes these reports as records of mathematical workshops and their open-problem discussions.
\item Supplied literature reference, labeled ``Recent subspace-design rank results'': \url{https://arxiv.org/abs/2404.06513}.
\item Supplied literature reference, labeled ``Equal-rank subspace designs'': \url{https://digitalcommons.mtu.edu/math-fp/97/}.
\item Literature assessment dated exactly \texttt{2026-08-22}: the supplied curation labels the problem open and states that recent work proves rank results for important families of subspace designs while continuing to describe the generalized Hamada conjecture as open; it also states that equal-rank nongeometric examples do not settle the conjecture. This is supplied dated curation derived from the OWR record, not an independently verified assessment. Its time-sensitive status and cited claims require independent verification.
\item The registry labels the statement \texttt{corrected\_verified} and says its mathematical content was preserved while notation, typography, duplicate clauses, or omitted definitions were repaired. These are supplied editorial claims; they do not resolve the formulation issues identified here or establish mathematical correctness.
\end{itemize}
\end{document}
